%!TEX root = ../main.tex

\addcontentsline{toc}{chapter}{付録}
\chapter*{付録}

\section*{基本料金を考慮し買電コストを最小化するケースにおける動的計画モデルによる定式化}
本研究の第4章で扱った動的計画モデルでは基本料金を考慮するケースである，ケース3を対象外とした．
本節ではケース3について動的計画モデルで定式化を行い，時間的計算量について考察する．

\subsection*{状態変数の拡張}
式\eqref{eq:minimize_lp_3}で表される，ケース3の目的関数には全期間を通じた最大需要電力 $s^{\rm BYMax}$ に依存する項が含まれている．
\begin{equation}
    \text{Minimize} \quad  \sum_{k=1}^{K}  ( p_k^{\rm{BY} } s_k^{\rm{BY}} +  p^{\rm{BYMax} } s^{\rm{BYMax}}  )．
 \label{eq:minimize_lp_3_1}
\end{equation}
ここで，$s^{\rm BYMax} $は30分の平均買電力量の最大値であるが，ここでは議論を単純化するために各期の買電力量の最大値とする．
すなわち，$s^{\rm BYMax}= \max (s_1^{\rm BY},s_2^{\rm BY}, \ldots, s_k^{\rm BY})$ とする．
通常の動的計画法において，各期の決定は「現在の状態」のみに依存するという性質が成立している必要がある．
しかし，基本料金は過去の買電の履歴に依存するため，第4章で定義した蓄電池残量の状態 $s_k$ だけでは，
その時点で基本料金が確定せず，最適な行動を決定できない．

したがって，ケース3を動的計画モデルで扱うためには，状態空間を拡張し，
第 $k$ 期までの最大需要電力を保持する新たな状態変数 $h_k$ を導入する必要がある．
拡張された状態 $\tilde{s}_k$ を以下のように定義する．
\begin{equation}
    \tilde{s}_k = (s_k, h_k)．
\end{equation}
ここで，$s_k \in \{0, \dots, B-1\}$ は蓄電池残量の状態，$h_k$ は第 $k$ 期までの最大買電力量を表す離散化された状態である．

\subsection*{状態遷移関数とBellman方程式}
拡張された状態における状態遷移について考える．
第4章では状態遷移関数 $T$ を $s_k = T(s_{k-1}, a_k)$ と定義したが，本付録では拡張された状態 $\tilde{s}_k$ に対する状態遷移関数 $\tilde{T}$ を導入する．
蓄電池残量の遷移は第4章と同様であり，最大買電力量の状態 $h_k$ はこれまでの最大値 $h_{k-1}$と現在の期の買電量のうち小さくない方に更新される．
したがって，状態遷移関数 $\tilde{T}$ は以下のように定義される．

\begin{equation}
    \tilde{s}_k = \tilde{T}(\tilde{s}_{k-1}, a_k) = (T(s_{k-1}, a_k), \max\{h_{k-1}, s_k^{\mathrm{BY}}(s_{k-1}, a_k)\})
\end{equation}

これに伴い，Bellman方程式は，拡張された状態 $\tilde{s}_{k-1} = (s_{k-1}, h_{k-1})$ を用いて次のように書き換えられる．
\begin{equation}
    F_{k-1}^{*}(\tilde{s}_{k-1}) = \min_{a_k \in A_{k}(s_{k-1})} \{ R_k(\tilde{s}_{k-1}, a_k) + F_{k}^{*}(T'(\tilde{s}_{k-1}, a_k)) \}
\end{equation}


\subsection*{計算量の増大に関する考察}
ここで，最大買電力量 $h_k$の離散化数を$H$と定義する．このとき，本モデルにおける時間的計算量は $O(B\{ \frac{(B-1)(\Delta s_k^{\mathrm{Max}}(s_{k-1}) - \Delta s_k^{\mathrm{Min}}(s_{k-1}) )}{\overline{b}} +1\} KH)$ となり，ケース1，2，4と比較すると約$H$倍計算に時間がかかることが想定される．
よって，本研究の動的計画モデルにおいてケース3を扱うことは，時間的計算量の観点から現実的ではない．